
\subsubsection{Self-Repair and Iterative Refinement}

Chen et al. (2025) conducted a study of self-repair across model scales, demonstrating that self-repair yields minimum +4.9\% improvement on HumanEval and MBPP benchmarks. Their work established that 8B+ parameter models can benefit from prompt-based self-repair without fine-tuning. However, their experiments used raw compiler output without investigating alternative formatting.

The Self-Refine framework \citep{madaan2023selfrefine} introduced a general paradigm for iterative LLM self-critique and revision, focusing on the iterative structure of refinement rather than the format of feedback signals.

InspectCoder (2025) compared static and dynamic analysis feedback for self-repair, finding that debugger-based dynamic feedback can complement static analysis. While this work varies the type of analysis, it does not systematically vary how static analysis errors are formatted.

\subsubsection{Compiler Feedback Integration}

CompCoder (2024) demonstrated improvements in compilation success (44\% $\rightarrow$ 89\%) by using compiler feedback as a training signal. This work validates that compiler information is valuable but treats the feedback format as fixed during inference.

CodeRL \citep{shojaee2023coderl} applied actor-critic reinforcement learning with unit test signals, optimizing for test-passing behavior. This approach requires fine-tuning and does not address feedback presentation during inference.

The survey by Zhang et al. (2025) on LLM-compiler integration catalogues various approaches to combining these technologies. All surveyed methods treat compiler output format as given rather than as a design variable.

\subsubsection{Type-Aware Code Generation}

TyFlow (2025) introduced type-guided program synthesis, demonstrating that type checker integration during generation improves correctness. This work addresses generation rather than repair.

ReCode (2025) combines retrieval-augmented generation with static analysis for code repair. Their fine-grained retrieval approach implicitly reformats error information by retrieving relevant examples, but does not isolate the effect of format from retrieved context.

\subsubsection{Positioning}

The works above demonstrate the value of compiler feedback, type constraints, and retrieval augmentation. However, none systematically study error message format as an independent variable while controlling for information content through a scrambled control condition.
